\chapter{“Sekunde Ishirini”}
\label{sec:1}
Afisa wa ujasusi, akieleza nafasi yake katika msururu wa kulenga shabaha uliojengwa kuzunguka orodha za watu wa kuuawa zilizotengenezwa na mashine, alimwambia mwandishi wa habari kwamba alitumia sekunde ishirini kwa kila jina, muda wa kutosha tu kuhakikisha kwamba shabaha alikuwa mwanamume: “Haichukui muda mrefu kutambua kama mtu ana sauti ya kiume au ya kike.” Aliongeza kwamba kama binadamu hakuongeza thamani yoyote, isipokuwa kuwa muhuri wa idhini. “Iliokoa muda mwingi” \autocite{abraham2024lavender}. Kwa maelezo yake, hiyo ndiyo maana halisi ya “binadamu ndani ya mzunguko” chini ya kasi ya operesheni.\footnote{Utafiti kuhusu upendeleo wa kuamini mashine unatabiri ukabidhishaji huu: mahitaji mengine yanaposhindania umakini wao, watu wanaopewa pendekezo la mashine hulikagua kidogo zaidi na kulitegemea zaidi \autocite{parasuraman2010complacency}.}

Tafsiri moja ni kwamba afisa huyo ni mhusika aliyepewa muda mdogo mno, na suluhisho ni kumtawala vizuri zaidi~— muda zaidi, mkaguzi wa pili, dhamiri ya ziada iliyowekwa ndani ya mzunguko. Kwa mtazamo huu wa \emph{utiishwaji wa kijamii}, mtu ni mhusika anayesimama nje ya mashine na kuiendesha kama chombo: dereva na gari, fundi na kereza. Mtu ana nafasi, jina, majukumu, na dhamiri inayoweza kuzungumzishwa. Mtazamo mbadala ni \emph{utumwa wa kimashine}, ambapo mtu ni kijenzi cha mashine, kilichounganishwa pamoja na vipande vyake vingine kuwa mfumo unaoviongoza vyote kwa pamoja; mashine inamtumia mtu kama mojawapo ya viungo vyake, na kile ambacho mtu huchangia ni kazi inayofanywa kwa kasi ambayo mfumo huiweka.\footnote{Jina ambalo Doctorow hulipa nafasi ya afisa huyo ni kentauri wa kinyume: si mtu anayesaidiwa na mashine bali kijenzi cha mashine kilichotengenezwa kwa mtu \autocite{doctorow2026reversecentaur}.} Mitazamo hii miwili haitengani; mtu anaweza kuwa na nafasi upande mmoja huku akiwa kipande upande mwingine \autocite[Uwanda 13, “The Apparatus of Capture”: utumwa upo pale binadamu wenyewe wanapokuwa vipande vinavyounda mashine ambayo wao huiunda miongoni mwao na pamoja na vitu vingine, chini ya udhibiti na uongozi wa umoja wa juu zaidi, na utiishwaji upo pale umoja huo wa juu zaidi unapomfanya binadamu kuwa mhusika aliyefungamanishwa na kitu kipya cha nje, ambacho kinaweza kuwa mnyama, chombo, au hata mashine]{deleuze1987plateaus}, na afisa wa sekunde ishirini alikuwa vyote viwili. Kama kipande, alikuwa mahali ambapo msururu ulitoa muhuri wake, na sekunde ishirini zilikuwa muda wa awamu moja ya kipande hicho. Kama mhusika, alizungumzishwa muda wote: alimweleza mwandishi wa habari mkinzano kati ya habari za vifo vya raia na shinikizo kazini la kuleta shabaha zaidi. Na hata mkinzano huo ulipoifikia dhamiri yake, aliendelea kukubali. Uliingia akilini na haukuamua chochote;\footnote{Au, angalau, haukuamua chochote kilichozuia Jeshi la Ulinzi la Israeli kuua watu 100,000 huko Gaza kati ya Oktoba 2023 na Oktoba 2025 \autocite{gomezugarte2025uncertainty,acosta2026gaza,economist2025gaza}. Na simulizi ya afisa huyo mahususi ilichapishwa Aprili 2024 na +972 Magazine. Hakuna anayejua kama uamuzi wa afisa huyo wa kuzungumza na mwandishi wa habari ulizuia kifo kimoja cha ziada, vifo 50,000 vya ziada, au hakuna kabisa. Na nini kingetokea kama afisa huyo angekataa baada ya wiki mbili katika nafasi hiyo~— tuseme, kwa dhana tu, Novemba 2023? Tunazungumzia afisa wa ujasusi anayetumia Lavender kukusanya orodha za watu wa kuuawa. Mazingira yake ya kazi ya kila siku pengine yalifanana zaidi na kazi ya ofisini ya kiraia kuliko filamu ya mapigano. Je, angepaswa kutafuta mikutano ya mara kwa mara na wakubwa wake, akipanda ngazi za uongozi, huku akikataa kuidhinisha shabaha za mashambulizi ya anga ambayo kwa busara yangetarajiwa kusababisha vifo vya raia? Huenda hilo ndilo lingekuwa jambo la kimaadili kufanya. Katika hali hiyo dhahania, je, kungekuwa na wafu wachache zaidi?} uamuzi kwamba sera hiyo haikuwa na uwiano ulikuja, alisema, baadaye alipotafakari.

Kwa hiyo kilichokosekana hakikuwa dhamiri; alikuwa nayo. Dhamiri hufanya kazi kwa kufanya uamuzi ugharimu kitu, lakini hapa gharama zilielekea upande mwingine: “Tulikuwa tukishinikizwa kila mara: ‘Tuleteeni shabaha zaidi.’ Kwa kweli walitufokea” \autocite{abraham2024lavender}. Alikuwa kijenzi chenye wasiwasi wa kimaadili, lakini kasi ilidumu: wasiwasi huo haukubadilisha chochote alichochangia, na akawa muhuri. Kijenzi chenye wasiwasi wa kimaadili kisichoendana na kasi kinakuwa “kizuizi cha kibinadamu” ambacho kamanda wa zamani wa kitengo cha Israeli cha ujasusi wa mawasiliano aliandika kitabu kuhusu kukiondoa \autocite{ys2021humanmachine}. Kilichomfanya afisa huyo kuwa kijenzi ni utaratibu, si mtazamo wake. Kijenzi kinaweza kufuata kwa kusita au kwa hamu: mtendaji asiyefikiri ambaye Arendt alimwona katika kesi ya Eichmann \autocite{arendt1963eichmann} alikuwa, kwa ushahidi wa yale Eichmann aliyoyasema Argentina miaka minne kabla, wakati hakuna kilichokuwa hatarini kwake, maigizo aliyoyafanya kwa ajili ya Yerusalemu \autocite{stangneth2014eichmann}. Kile ambacho msururu ulimruhusu afisa huyo kuchangia kilikuwa kazi moja isiyobadilika kwa kasi isiyobadilika: aliona kila kesi kama msururu ulivyoieleza, moja baada ya nyingine, katika muda ambao msururu ulimpa. Kama kuna kitu chochote kinachoweza kusimama mahali hapo bila kuwa kijenzi ndilo swali la Sehemu ya IV.

Uandishi wa habari za uchunguzi ulieleza mifumo mitatu iliyofanya kazi pamoja huko Gaza: mmoja ukizalisha shabaha za majengo na miundombinu, huku vyanzo vikisema mapendekezo mengi yalikuwa makazi ya watu binafsi; wa pili ukiwapa watu mmoja mmoja alama kulingana na makadirio ya uhusiano wao na wanamgambo, ukitoa orodha zinazosemekana kuwa ziliwawekea alama watu wapatao 37,000, ambazo ndani ya jeshi zilieleweka kuwa zinakosea takriban mara moja kati ya kumi\footnote{Jeshi la Israeli lilipinga simulizi hiyo. Simulizi za Minab na Maven hapa chini zinategemea uandishi wa habari za uchunguzi ambao pande husika zinaupinga pia. Ninazichukua kama zilivyoripotiwa.}; na wa tatu ukifuatilia wakati watu waliowekewa alama walirudi nyumbani, ili mashambulizi ya anga yapangwe wakati walipokuwa na familia zao.

\runin{Maven: toleo la Kimarekani, sasa likiongezwa kasi na akili bandia zalishi}

Shambulio la bomu dhidi ya shule ya Minab liliua watu 156, 120 kati yao wakiwa watoto wa shule \autocite[aya~92]{UNFFMIran2026}.\footnote{Idadi zinatofautiana kulingana na chanzo na tarehe, kutoka 156 (idadi ya mwisho ya mwendesha mashtaka wa Minab, ambayo pia imetolewa na Amnesty International) hadi 168 (meya wa mji huo, katika siku za kwanza). Ujumbe wa Umoja wa Mataifa unatoa takwimu za mwendesha mashtaka na pia unaripoti taarifa huru za vifo 157, 123 kati yao wakiwa watoto wenye umri wa miaka 13 au chini \autocite[aya~92–93]{UNFFMIran2026}.} Kazi ya kiraia ya shule hiyo ilionekana katika picha za setilaiti zinazopatikana hadharani, na usimbaji wa shabaha bado ulikuwa na alama iliyorithiwa kutoka kambi jirani ya Jeshi la Walinzi wa Mapinduzi. Ujumbe wa Umoja wa Mataifa wa kutafuta ukweli uliripoti kwamba kuna sababu za msingi za kuamini kuwa shambulio hilo lilikuwa uhalifu wa kivita wa kukifanya kitu cha kiraia kuwa lengo la shambulio \autocite{ohchr2026minabpress}, na kwamba shule hiyo ilitambulika waziwazi bila taarifa yoyote inayoonyesha matumizi ya kijeshi \autocite[aya~81–86, 124–126]{UNFFMIran2026}. Ujumbe huo uligundua kwamba Marekani ilishindwa kutimiza wajibu wake wa kufanya kila linalowezekana kuthibitisha shabaha, ukihesabu usasishaji wa data za kulenga shabaha miongoni mwa tahadhari zinazowezekana, na ukakuweka kushindwa huko juu ya uzembe \autocite[aya~124–125]{UNFFMIran2026}. Amnesty International ilifikia hitimisho linalofanana, kwa lugha inayofuata wajibu wa kuthibitisha katika Ibara ya 57(2)(a)(i) ya Itifaki ya Ziada ya Kwanza: kwamba majeshi huenda yalitegemea taarifa za kijasusi zilizopitwa na wakati na yakashindwa kufanya kila linalowezekana kuthibitisha shabaha \autocite{Amnesty2026Minab}.

Ilitokea asubuhi ya tarehe 28 Februari, siku ya kwanza ya kampeni ya Iran \autocite{Amnesty2026Minab,crs2026iranwar}. Siku moja kabla, Waziri alikuwa ametangaza kwamba muuzaji (Anthropic) ambaye modeli yake (Claude) ilikuwa ndani ya jukwaa la kulenga shabaha (Maven Smart System) angetangazwa kuwa hatari kwa mnyororo wa ugavi kwa kukataa kutoa leseni ya Claude kwa matumizi yote halali \autocite{npr2026openai}. Maven Smart System ilikuwa imeirudisha shule hiyo kama shabaha iliyopendekezwa kwa siku ya kwanza \autocite{livingstone2026maven,gizmodo2026minab}. Hakuna simulizi iliyochapishwa inayoeleza Claude ilikuwa na nafasi gani katika pendekezo hilo, na mtu mmoja anayeifahamu kazi ya muuzaji huyo aliiambia NBC News kwamba modeli hiyo haikutoa moja kwa moja ushauri wa kulenga shabaha \autocite{nbc2026targeting}.

Ukaguzi wa ndani wa Pentagon, ambao matokeo yake kamili bado yanafichwa kwa usiri mkubwa, unaripotiwa kugundua kwamba maandalizi ya shabaha yaliyochukua saa kadhaa katika migogoro ya awali yalibanwa hadi dakika, kwamba wafanyakazi wa kupunguza madhara kwa raia walipungua kwa takriban asilimia tisini hadi chini ya watu ishirini, kwamba timu ya Central Command ilipungua kutoka watu kumi hadi mmoja, na kwamba hakuna mwanachama wake aliyekagua eneo la Minab~— matokeo ambayo maafisa waliyahusisha na upunguzaji huo na pia na uamuzi wa kutoihusisha timu hiyo katika mipango na operesheni \autocite{Bloomberg2026KillChain}. Uchunguzi rasmi ulikuwa umefunguliwa Machi, baada ya tathmini ya awali kubaini kwamba huenda Marekani ndiyo iliyohusika \autocite{NPR2026MinabProbe}.

Katika kampeni ya Iran ya 2026, jukwaa la kulenga shabaha lililojengwa na mkandarasi, likiwa na modeli ya matumizi ya jumla ya muuzaji wa kibiashara ndani yake, na likitumika kupanga shabaha kwa umuhimu wa kimkakati na kuandaa rasimu ya uhalali wa kisheria kwa kila shambulio \autocite{livingstone2026maven}, lilisaidia kushambulia zaidi ya shabaha elfu moja katika saa ishirini na nne za kwanza, takriban mara kumi ya kilichowezekana kabla yake \autocite{csis2026maven}, na shabaha elfu kumi na tatu kufikia siku ya thelathini na nane \autocite{breakingdefense2026maven}. Pentagon haijachapisha muda wa ukaguzi kwa kila shabaha kwa jukwaa hilo.\footnote{Maven ni jina la mpango mmoja linalobeba zaidi ya mfumo mmoja: mpango wa Pentagon ambao bado unaitwa Maven, na Maven Smart System, jukwaa la programu lililotokana nao \autocite{csis2026maven}. Mpango huo ulianza mwaka 2017 kama mradi wa kuweka lebo kwenye vitu katika video ambazo ndege zisizo na rubani zilikuwa zikikusanya wakati wa kampeni dhidi ya Dola la Kiislamu \autocite{weisgerber2017maven,fang2019maven}. Mfumo wa 2026 unaunganisha mifumo ya kijeshi na vyanzo vya ujasusi katika kiolesura kimoja kinachotumika kutafuta na kushambulia shabaha na kuongoza operesheni \autocite{defensescoop2026transition}. Jukwaa hilo ni la Palantir: afisa mkuu wa masuala ya kidijitali na akili bandia wa Pentagon analiita “bidhaa ya Palantir ya programu kama huduma” \autocite{defensescoop2026transition}. Msimamo wa hadharani wa Palantir, ulionukuliwa kupitia mtu mwingine, ni kwamba kuamua uhai au kifo si jukumu lake na kwamba jukumu hilo daima hubaki kwa taasisi ya kijeshi \autocite{lee2026palantir}, na kampuni hiyo inashikilia kwamba taarifa za kijasusi za msingi zilikuwa jukumu la serikali na kwamba hakuna ushahidi kwamba Maven ilifanya kazi vibaya au ilisababisha shambulio la Minab \autocite{Bloomberg2026KillChain}. Jinsi mpango huo ulivyohama kutoka ofisi moja na mkandarasi mmoja hadi mwingine imeelezwa katika Kiambatisho~\ref{sec:C.1}.}

Watu waliokuwa kwenye orodha za mashambulizi ya anga za Israeli na Marekani hawangeweza kupinga kuwekwa kwenye orodha hizo~— misururu hiyo haikujengwa kupokea maoni ya aina hiyo. Andry José Hernández Romero alinaswa katika msururu tofauti, na yeye angeweza kupinga. Msanii wa vipodozi kutoka Venezuela, alijitokeza katika kituo cha kuingia cha San Ysidro mnamo Agosti 2024 kuomba hifadhi, akafaulu mahojiano ya kuthibitisha hofu ya kweli ya kuteswa, na alikuwa na wakili, ombi lililokuwa likisubiri uamuzi, na tarehe ya kusikilizwa \autocite{LULAC2025HernandezRomero,Vega2025HernandezRomero}. Mfanyakazi wa mkandarasi aliyekuwa akimchunguza alitafsiri tatoo zake mbili, taji juu ya maneno ya “mama” na “baba”, kama alama za genge \autocite{Palomera2025Blade}, na Machi 2025 alisafirishwa kwa ndege hadi gereza moja huko El Salvador kabla ya kusikilizwa kwa kesi yake, jioni ileile ambayo jaji wa shirikisho aliamuru ndege hizo zirudishwe \autocite{jgg2025contempt}. Alikuwa na kila kitu ambacho mtu anahitaji ili kupinga isipokuwa muda. Tutarudi kwenye kesi yake kwa kina zaidi tutakapojadili msururu wa Marekani wa utekaji nyara na uhamishaji wa kulazimishwa.

\textbf{Nani alikataa.} Kumbukumbu ya sekunde ishirini ina nusu ya pili, kumbukumbu ya nani alistahimili shinikizo la aina hiyo, na inasema kwamba mpangilio wa chumba una umuhimu mkubwa. Katika hali ya msingi ya Milgram, karibu theluthi moja ya washiriki waliacha kabla ya mshtuko wa juu kabisa wa umeme. Wenzao wawili walipokaidi, ni mmoja tu kati ya kumi aliyefika mwisho, na wasimamizi wawili wa jaribio walipotofautiana, wachache zaidi sana waliendelea kuliko katika hali ya msingi \autocite{milgram1974obedience}. Takriban robo ya washiriki wa Asch hawakuwahi kukubaliana na walio wengi waliokosea \autocite{asch1956studies}. Browning anakadiria kwamba kati ya sehemu moja ya kumi na sehemu moja ya tano ya wanaume wa Kikosi cha Polisi wa Akiba 101 walikataa au walikwepa ufyatuaji risasi huko Józefów, mauaji ya halaiki ya kwanza ya kikosi hicho \autocite{browning1992ordinary}. Uwiano huo ulitofautiana kutoka karibu hakuna hadi karibu wote kulingana na mazingira. Watu waliostahimili walikuwa na ushahidi kutoka nje ya njia rasmi: mwenzao aliyekaidi, mamlaka ya pili iliyotofautiana, na, miongoni mwa waokoaji ambao akina Oliner waliwachunguza, mahusiano yaliyovuka mstari ambao mateso yalichora \autocite{oliner1988altruistic}. Au walishikilia msimamo wao kwa usahihi ambao shinikizo halingeweza kuushinda, kama misimamo ya kimaadili isiyojadilika na maadili matakatifu yanavyoshikiliwa \autocite{skitka2009limits,ginges2007sacred}. Au walikichukulia chanzo kuwa na maslahi yake binafsi. Na makundi yalipinga mahali ambapo watu waliokuwa peke yao mara chache walifanya hivyo \autocite{gamson1982encounters}.

Chumba kinaweza kupangwa vilevile kwa ajili ya upande mwingine. Ujerumani wa enzi ya Weimar ulikuwa na maisha tele ya vyama na vikundi vya kijamii, ambayo huenda yalisaidia kuiangusha jamhuri: kadiri vyama vya siasa na serikali vilivyoshindwa kujibu malalamiko ya watu, watu waliweka nguvu zao katika vilabu na jumuiya zilizopangwa kwa misingi ya makundi, na Wanazi waliajiri wanachama kupitia mitandao ambayo jumuiya hizo zilikuwa zimeshaijenga \autocite{berman1997weimar}. Miji yenye kwaya, vilabu vya michezo na vyama vya mazoezi ya viungo vingi zaidi iliona watu wakijiunga na chama hicho kwa kasi zaidi, lakini pale tu serikali ilipokuwa haina utulivu; huko Prussia, ambako serikali ilikuwa tulivu, uhusiano huo ulikuwa mdogo na haukuwa na umuhimu wa kitakwimu \autocite{satyanath2017bowling}. Kwa hiyo kile ambacho waokoaji walikuwa nacho hakikuwa tu kuwa na watu wengi zaidi karibu nao. Mahusiano yao yalivuka mstari ambao mateso yalichora. Ashutosh Varshney aligundua jambo hilohilo katika miji ya India: ghasia kati ya Wahindu na Waislamu zilikuwa nadra pale ambapo vyama vya kiraia viliziweka jamii hizo mbili pamoja katika mashirika yaleyale, na za kawaida pale ambapo havikufanya hivyo \autocite{varshney2002ethnic}. Mahusiano yanayowezesha kukataa ndiyo yale ambayo vuguvugu la kifashisti lingelazimika kuyakata.

Kila kinachofuata kinahusu kupanga chumba.\footnote{Chumba ni mandhari ya jukwaa, na kama mwalimu wangu wa maigizo angesema, mandhari huongoza onyesho. Nani mwingine yumo chumbani? Je, mtu anaweza kuondoka na akanusurika? Je, kuna yeyote angesikia mtu akikataa?} Kupanga ni neno la kutengeneza kitu, na Hannah Arendt alionya dhidi ya kuichukulia siasa kama kitu kinachotengenezwa. Kwa Wagiriki, aliandika, sheria ilikuwa ukuta kihalisi, ambao bila huo kungeweza kuwa na kijiji lakini si mji \autocite{arendt1958human}: mtunga sheria alijenga ukuta, na kilichotokea ndani yake hakikuwa chake kukitengeneza. Hatua za msingi ninazozitetea ni kuta kwa maana hiyo. Zinamwezesha anayekataa kunusurika na kusikika; hazikizalishi kukataa, na mamlaka, kwa mujibu wa Arendt, yapo tu kati ya watu wanaotenda pamoja na hutoweka wanapotawanyika, jambo ambalo kumbukumbu ya nani alikataa inaonyesha kuhusu kukataa. Vyombo vya Sehemu ya II na III vinafanya kukataa kuhesabike. Taasisi za Sehemu ya V zinazalisha wakataaji. Mashine inayoweza kukataa ingehitaji kile ambacho wale tisa kati ya kumi waliokuwa na mwenzao kando yao walikuwa nacho: mwenzake mwenye historia tofauti kando yake, mahali pa kwenda, na mtu ambaye angeisikia.

Hicho pia ndicho ambacho mkataaji~— ikiwemo mkataaji wa kimashine~— katika msururu angeweza kuwa kwa afisa huyo. Mashine iliyokataa ingekuwa kwenye skrini wakati kesi inamfikia: mwenzake anayekaidi, ikiwa athari hiyo itastahimili kuhama kutoka chumbani hadi kwenye skrini, jambo ambalo halijajaribiwa na ambalo upendeleo wa kuamini mashine unatabiri halitatokea. Mashine iliyowasukuma wakaguzi kwa sababu tu waliinyenyekea ingewafanya kuwa mihuri ya upande wa pili. Je, afisa mwenye sekunde ishirini na mashine inayokataa kando yake ni afisa yuleyule kama yule wa kwenye ripoti?
